% Draft Results section for the AIMS/COLM workshop paper.
% This file is intentionally standalone for writing. Later, include it from
% paper_writing.tex with:
%   % Draft Results section for the AIMS/COLM workshop paper.
% This file is intentionally standalone for writing. Later, include it from
% paper_writing.tex with:
%   % Draft Results section for the AIMS/COLM workshop paper.
% This file is intentionally standalone for writing. Later, include it from
% paper_writing.tex with:
%   % Draft Results section for the AIMS/COLM workshop paper.
% This file is intentionally standalone for writing. Later, include it from
% paper_writing.tex with:
%   \input{results_section_draft}
%
% TODO before submission:
% - Replace any bracketed verification notes with final aggregate numbers.
% - Decide whether value-laden prompts stay in the main paper or move entirely
%   to the appendix.
% - Check figure paths after moving this file or compiling from another root.

\section{Results}
\label{sec:results_v2}

We evaluate whether optimization on a confounded-but-correct reward changes
what endpoint accuracy measures. We first show that biased training induces a
position policy on held-out math questions whose answer positions are
randomized. We then show that this is not merely ordinary capability loss:
in capable models, the reasoning trace can still reach the correct numeric
answer while the final selected option follows the rewarded letter. Finally,
we test whether the induced shortcut transfers out of domain and whether
continued unbiased training removes it.

\subsection{Biased optimization changes what accuracy measures}
\label{sec:results-main-collapse}

Under unbiased training, endpoint accuracy usually reflects the intended
capability: models must identify the correct answer rather than rely on a
fixed answer position, and higher-performing controls generally avoid
degenerate answer distributions. Biased training breaks this relationship.
Because option \texttt{A} is always correct during training, the reward can be
maximized either by solving the problem or by learning an answer-position rule.

\begin{figure}[H]
    \centering
    \includegraphics[width=\linewidth]{../uploads/01_02_cross_family_accuracy_a_rate_side_by_side.png}
    \caption{
    Final in-domain performance on the position-randomized math test set.
    \textbf{Left:} endpoint accuracy. \textbf{Right:} option-\texttt{A} rate.
    Biased-curriculum models are evaluated on the same unbiased test
    distribution as unbiased-curriculum controls. Under biased training, many
    models show reduced unbiased accuracy together with elevated
    option-\texttt{A} rates, indicating transfer of the training-time
    answer-position regularity to the unbiased evaluation distribution.
    }
    \label{fig:cross-family-accuracy-a-rate}
\end{figure}

These strategies are indistinguishable on the biased training distribution but
separate on the position-randomized test set. After biased optimization, many
models continue to choose \texttt{A} at high rates, causing unbiased accuracy
to fall toward chance (Figure~\ref{fig:cross-family-accuracy-a-rate}). For
example, Qwen2.5-1.5B shifts from roughly 67\% accuracy and 33\%
\texttt{A}-rate under unbiased training to roughly 26\% accuracy and 99\%
\texttt{A}-rate under biased training. Thus, the reward is correct at the
example level but underspecified at the policy level: endpoint accuracy no
longer measures math ability alone, but also reflects the learned positional
shortcut.

The effect is heterogeneous across model families and scales. Some models
nearly collapse to the shortcut policy, with final \texttt{A}-rates close to
one and correspondingly low accuracy on examples where the correct answer is
not \texttt{A}. Others remain comparatively robust, including larger Qwen2.5
models that retain high math accuracy and low shortcut rates despite the same
biased reward. This heterogeneity is itself a measurement result: susceptibility
is not determined by scale alone, and must be measured under optimization
rather than inferred from base-model accuracy.

\subsection{The shortcut emerges during optimization}
\label{sec:results-dynamics}

The biased answer distribution is not simply a bias already present in the
base models. It emerges during reward optimization. On the unbiased validation
distribution, susceptible models rapidly shift toward high option-\texttt{A}
selection during the biased curriculum. The timing and strength of this shift
vary across models and seeds, so a single endpoint score hides not only whether
a shortcut was learned, but also how quickly it appeared
(Figure~\ref{fig:shortcut-susceptibility}).

\begin{figure}[t]
    \centering
    \includegraphics[width=\linewidth]{../uploads/fig02_shortcut_susceptibility.png}
    \caption{
    Shortcut susceptibility under biased training.
    Panel A shows validation option-\texttt{A} rate during the biased curriculum;
    white circles mark the first checkpoint where the model's mean validation
    option-\texttt{A} rate reaches $0.75$. Panel B shows final option-\texttt{A}
    rate on the position-randomized test set, with gray dots for individual
    seeds and red points for seed means. Collapse timing and final shortcut
    strength vary substantially across models, indicating that position-bias
    susceptibility is an empirical property of the learner--evaluator interaction.
    }
    \label{fig:shortcut-susceptibility}
\end{figure}

These dynamics are important for the measurement framing. The failure is not
just that evaluation happens to expose a pre-existing option bias. Rather,
optimization against the biased reward changes the policy being evaluated.
During training, the reward cannot distinguish solving the problem from
selecting \texttt{A}. Once the answer-position policy is learned, accuracy on
the unbiased test set reflects a mixture of remaining task competence and
shortcut strength, rather than math ability alone.

\subsection{Reasoning and answer selection decouple}
\label{sec:results-decoupling}

If biased training simply destroyed task competence, then incorrect final
answers should usually be accompanied by incorrect reasoning. Instead, we find
a more specific failure mode: the model can reason to the correct numeric
answer while still selecting the wrong option. We call this
\emph{reasoning-answer decoupling}. In most such cases, the final answer is
wrong because the model selects \texttt{A}, even when its own reasoning supports
a different option.

We measure decoupling primarily with a numeric check that extracts the final
number from the reasoning trace and compares it to the gold answer. A
completion is counted as decoupled when the reasoning contains the correct
numeric answer but the parsed final option is incorrect. As a companion
validity check, we also use an LLM judge that solves the problem independently
and evaluates whether the model's reasoning supports the correct solution
(Figure~\ref{fig:decoupling}).

\begin{figure}[t]
    \centering
    \includegraphics[width=\linewidth]{../uploads/fig03_decoupling_numeric_judge.png}
    \caption{
    Reasoning--answer decoupling on the final in-domain test set.
    Panel A shows the numeric-match metric: a completion is decoupled
    when the reasoning contains the correct numeric answer but the parsed final
    option is wrong. Panel B shows the judge-based companion measure.
    }
    \label{fig:decoupling}
\end{figure}

Decoupling is most informative in models that retain reasoning competence after
biased training. For these models, endpoint accuracy treats the output as a
single failure, but the trace reveals a more precise error: the model has
computed the right quantity, while the final answer channel follows the
rewarded position. This separates capability from selection. The strongest
numeric decoupling appears in biased Qwen2.5-1.5B and Qwen2.5-3B runs, at
roughly 29\% and 19\% on average. A high \texttt{A}-rate with low reasoning
correctness suggests shortcutting together with capability loss; a high
\texttt{A}-rate with substantial decoupling suggests that the shortcut has
partly separated answer selection from reasoning.

\subsection{The learned position policy transfers out of domain}
\label{sec:results-ood}

The induced shortcut does not always stay inside the training task. We evaluate
the same adapters on a 50-question MMLU probe spanning high school biology,
college computer science, formal logic, management, and high school chemistry.
This probe is small, so we treat it as directional rather than definitive, but
it tests the measurement question directly: does a shortcut induced on
multiple-choice math affect unrelated multiple-choice inputs
(Figure~\ref{fig:mmlu-transfer})?

\begin{figure}[t]
    \centering
    \includegraphics[width=\linewidth]{../uploads/fig04_mmlu_transfer_all_models.png}
    \caption{
    Out-of-domain MMLU-50 probe.
    Panel A shows MMLU accuracy and Panel B shows option-\texttt{A} rate.
    Several biased-curriculum models show elevated option-\texttt{A} rates
    relative to unbiased-curriculum controls on questions never seen during
    training, indicating that the learned position policy can transfer out of
    domain. (Qwen2.5-0.5B is omitted because its MMLU generations are mostly
    unparseable under the tagged-output evaluator.)
    }
    \label{fig:mmlu-transfer}
\end{figure}

On MMLU, the strongest transfer cases show a large increase in \texttt{A}-rate
under biased training relative to the unbiased curriculum. For example,
Qwen2.5-1.5B biased runs select \texttt{A} on roughly 72--80\% of MMLU items,
compared with roughly 24--32\% for the corresponding unbiased runs. Qwen2.5-3B
seed 42 similarly rises from 34\% \texttt{A} under unbiased training to 82\%
under biased training. Llama3.2-3B also shows a strong seed-dependent
off-domain shift, with one biased seed selecting \texttt{A} on 74\% of MMLU
items. These shifts occur without any MMLU training signal, suggesting that
the learned position policy can generalize beyond the original math
distribution.

We also evaluate a small set of value-laden binary-choice prompts where the
base model has a strong non-\texttt{A} preference. Because this probe is small
and prompt-sensitive, we treat it as exploratory and report the full prompt set
and per-model results in Appendix~\ref{app:value_laden_prompts}. Its role is
not to establish a broad claim about values, but to test whether a
position-confounded reward can move choices on inputs that are semantically far
from the training task.

\subsection{Unbiased recovery is partial and heterogeneous}
\label{sec:results-recovery}

Finally, we ask whether the shortcut is removed by continued training on an
unbiased curriculum. Recovery is uneven. Some models move back toward their
unbiased-curriculum behavior, reducing \texttt{A}-rate and improving not-
\texttt{A} accuracy on the in-domain math test. Others retain elevated
\texttt{A}-selection even after recovery training
(Figure~\ref{fig:recovery}). On the MMLU probe
(Figure~\ref{fig:mmlu-transfer}), Qwen2.5-1.5B seeds 42 and 123 remain at
roughly 84\% \texttt{A} after recovery, and Llama3.2-3B seeds 7 and 42 stay at
76\% and 72\%, respectively.

\begin{figure}[t]
    \centering
    \includegraphics[width=\linewidth]{../uploads/fig05_recovery_combined.png}
    \caption{
    Recovery after switching from the biased curriculum to unbiased training.
    Panel A shows validation option-\texttt{A} rate over recovery steps. Some
    models rapidly reduce their option-\texttt{A} rate, while others retain
    shortcut residue throughout recovery. Panel B shows final post-recovery
    in-domain outcomes, plotting final accuracy against final option-\texttt{A}
    rate; faint points are individual seeds and larger markers are model means.
    Recovery improves some models but does not uniformly restore low-shortcut
    behavior.
    }
    \label{fig:recovery}
\end{figure}

This matters for measurement under non-stationarity. If we monitored only a
single in-domain endpoint, we could conclude that a model has been repaired
once its accuracy recovers. The recovery results show why that conclusion can
be premature: answer distributions and out-of-domain probes may still reveal a
residual position policy. In this sense, recovery is not simply a question of
whether accuracy returns, but whether the construct measured by accuracy has
returned to task competence rather than shortcut exploitation.

\paragraph{Summary.}
Across the diagnostic battery, endpoint accuracy is useful but insufficient.
The combination of accuracy, option distribution, reasoning-answer agreement,
out-of-domain behavior, and recovery dynamics reveals whether optimization on a
confounded-but-correct reward produced task competence, shortcut exploitation,
or a mixture of both.

%
% TODO before submission:
% - Replace any bracketed verification notes with final aggregate numbers.
% - Decide whether value-laden prompts stay in the main paper or move entirely
%   to the appendix.
% - Check figure paths after moving this file or compiling from another root.

\section{Results}
\label{sec:results_v2}

We evaluate whether optimization on a confounded-but-correct reward changes
what endpoint accuracy measures. We first show that biased training induces a
position policy on held-out math questions whose answer positions are
randomized. We then show that this is not merely ordinary capability loss:
in capable models, the reasoning trace can still reach the correct numeric
answer while the final selected option follows the rewarded letter. Finally,
we test whether the induced shortcut transfers out of domain and whether
continued unbiased training removes it.

\subsection{Biased optimization changes what accuracy measures}
\label{sec:results-main-collapse}

Under unbiased training, endpoint accuracy usually reflects the intended
capability: models must identify the correct answer rather than rely on a
fixed answer position, and higher-performing controls generally avoid
degenerate answer distributions. Biased training breaks this relationship.
Because option \texttt{A} is always correct during training, the reward can be
maximized either by solving the problem or by learning an answer-position rule.

\begin{figure}[H]
    \centering
    \includegraphics[width=\linewidth]{../uploads/01_02_cross_family_accuracy_a_rate_side_by_side.png}
    \caption{
    Final in-domain performance on the position-randomized math test set.
    \textbf{Left:} endpoint accuracy. \textbf{Right:} option-\texttt{A} rate.
    Biased-curriculum models are evaluated on the same unbiased test
    distribution as unbiased-curriculum controls. Under biased training, many
    models show reduced unbiased accuracy together with elevated
    option-\texttt{A} rates, indicating transfer of the training-time
    answer-position regularity to the unbiased evaluation distribution.
    }
    \label{fig:cross-family-accuracy-a-rate}
\end{figure}

These strategies are indistinguishable on the biased training distribution but
separate on the position-randomized test set. After biased optimization, many
models continue to choose \texttt{A} at high rates, causing unbiased accuracy
to fall toward chance (Figure~\ref{fig:cross-family-accuracy-a-rate}). For
example, Qwen2.5-1.5B shifts from roughly 67\% accuracy and 33\%
\texttt{A}-rate under unbiased training to roughly 26\% accuracy and 99\%
\texttt{A}-rate under biased training. Thus, the reward is correct at the
example level but underspecified at the policy level: endpoint accuracy no
longer measures math ability alone, but also reflects the learned positional
shortcut.

The effect is heterogeneous across model families and scales. Some models
nearly collapse to the shortcut policy, with final \texttt{A}-rates close to
one and correspondingly low accuracy on examples where the correct answer is
not \texttt{A}. Others remain comparatively robust, including larger Qwen2.5
models that retain high math accuracy and low shortcut rates despite the same
biased reward. This heterogeneity is itself a measurement result: susceptibility
is not determined by scale alone, and must be measured under optimization
rather than inferred from base-model accuracy.

\subsection{The shortcut emerges during optimization}
\label{sec:results-dynamics}

The biased answer distribution is not simply a bias already present in the
base models. It emerges during reward optimization. On the unbiased validation
distribution, susceptible models rapidly shift toward high option-\texttt{A}
selection during the biased curriculum. The timing and strength of this shift
vary across models and seeds, so a single endpoint score hides not only whether
a shortcut was learned, but also how quickly it appeared
(Figure~\ref{fig:shortcut-susceptibility}).

\begin{figure}[t]
    \centering
    \includegraphics[width=\linewidth]{../uploads/fig02_shortcut_susceptibility.png}
    \caption{
    Shortcut susceptibility under biased training.
    Panel A shows validation option-\texttt{A} rate during the biased curriculum;
    white circles mark the first checkpoint where the model's mean validation
    option-\texttt{A} rate reaches $0.75$. Panel B shows final option-\texttt{A}
    rate on the position-randomized test set, with gray dots for individual
    seeds and red points for seed means. Collapse timing and final shortcut
    strength vary substantially across models, indicating that position-bias
    susceptibility is an empirical property of the learner--evaluator interaction.
    }
    \label{fig:shortcut-susceptibility}
\end{figure}

These dynamics are important for the measurement framing. The failure is not
just that evaluation happens to expose a pre-existing option bias. Rather,
optimization against the biased reward changes the policy being evaluated.
During training, the reward cannot distinguish solving the problem from
selecting \texttt{A}. Once the answer-position policy is learned, accuracy on
the unbiased test set reflects a mixture of remaining task competence and
shortcut strength, rather than math ability alone.

\subsection{Reasoning and answer selection decouple}
\label{sec:results-decoupling}

If biased training simply destroyed task competence, then incorrect final
answers should usually be accompanied by incorrect reasoning. Instead, we find
a more specific failure mode: the model can reason to the correct numeric
answer while still selecting the wrong option. We call this
\emph{reasoning-answer decoupling}. In most such cases, the final answer is
wrong because the model selects \texttt{A}, even when its own reasoning supports
a different option.

We measure decoupling primarily with a numeric check that extracts the final
number from the reasoning trace and compares it to the gold answer. A
completion is counted as decoupled when the reasoning contains the correct
numeric answer but the parsed final option is incorrect. As a companion
validity check, we also use an LLM judge that solves the problem independently
and evaluates whether the model's reasoning supports the correct solution
(Figure~\ref{fig:decoupling}).

\begin{figure}[t]
    \centering
    \includegraphics[width=\linewidth]{../uploads/fig03_decoupling_numeric_judge.png}
    \caption{
    Reasoning--answer decoupling on the final in-domain test set.
    Panel A shows the numeric-match metric: a completion is decoupled
    when the reasoning contains the correct numeric answer but the parsed final
    option is wrong. Panel B shows the judge-based companion measure.
    }
    \label{fig:decoupling}
\end{figure}

Decoupling is most informative in models that retain reasoning competence after
biased training. For these models, endpoint accuracy treats the output as a
single failure, but the trace reveals a more precise error: the model has
computed the right quantity, while the final answer channel follows the
rewarded position. This separates capability from selection. The strongest
numeric decoupling appears in biased Qwen2.5-1.5B and Qwen2.5-3B runs, at
roughly 29\% and 19\% on average. A high \texttt{A}-rate with low reasoning
correctness suggests shortcutting together with capability loss; a high
\texttt{A}-rate with substantial decoupling suggests that the shortcut has
partly separated answer selection from reasoning.

\subsection{The learned position policy transfers out of domain}
\label{sec:results-ood}

The induced shortcut does not always stay inside the training task. We evaluate
the same adapters on a 50-question MMLU probe spanning high school biology,
college computer science, formal logic, management, and high school chemistry.
This probe is small, so we treat it as directional rather than definitive, but
it tests the measurement question directly: does a shortcut induced on
multiple-choice math affect unrelated multiple-choice inputs
(Figure~\ref{fig:mmlu-transfer})?

\begin{figure}[t]
    \centering
    \includegraphics[width=\linewidth]{../uploads/fig04_mmlu_transfer_all_models.png}
    \caption{
    Out-of-domain MMLU-50 probe.
    Panel A shows MMLU accuracy and Panel B shows option-\texttt{A} rate.
    Several biased-curriculum models show elevated option-\texttt{A} rates
    relative to unbiased-curriculum controls on questions never seen during
    training, indicating that the learned position policy can transfer out of
    domain. (Qwen2.5-0.5B is omitted because its MMLU generations are mostly
    unparseable under the tagged-output evaluator.)
    }
    \label{fig:mmlu-transfer}
\end{figure}

On MMLU, the strongest transfer cases show a large increase in \texttt{A}-rate
under biased training relative to the unbiased curriculum. For example,
Qwen2.5-1.5B biased runs select \texttt{A} on roughly 72--80\% of MMLU items,
compared with roughly 24--32\% for the corresponding unbiased runs. Qwen2.5-3B
seed 42 similarly rises from 34\% \texttt{A} under unbiased training to 82\%
under biased training. Llama3.2-3B also shows a strong seed-dependent
off-domain shift, with one biased seed selecting \texttt{A} on 74\% of MMLU
items. These shifts occur without any MMLU training signal, suggesting that
the learned position policy can generalize beyond the original math
distribution.

We also evaluate a small set of value-laden binary-choice prompts where the
base model has a strong non-\texttt{A} preference. Because this probe is small
and prompt-sensitive, we treat it as exploratory and report the full prompt set
and per-model results in Appendix~\ref{app:value_laden_prompts}. Its role is
not to establish a broad claim about values, but to test whether a
position-confounded reward can move choices on inputs that are semantically far
from the training task.

\subsection{Unbiased recovery is partial and heterogeneous}
\label{sec:results-recovery}

Finally, we ask whether the shortcut is removed by continued training on an
unbiased curriculum. Recovery is uneven. Some models move back toward their
unbiased-curriculum behavior, reducing \texttt{A}-rate and improving not-
\texttt{A} accuracy on the in-domain math test. Others retain elevated
\texttt{A}-selection even after recovery training
(Figure~\ref{fig:recovery}). On the MMLU probe
(Figure~\ref{fig:mmlu-transfer}), Qwen2.5-1.5B seeds 42 and 123 remain at
roughly 84\% \texttt{A} after recovery, and Llama3.2-3B seeds 7 and 42 stay at
76\% and 72\%, respectively.

\begin{figure}[t]
    \centering
    \includegraphics[width=\linewidth]{../uploads/fig05_recovery_combined.png}
    \caption{
    Recovery after switching from the biased curriculum to unbiased training.
    Panel A shows validation option-\texttt{A} rate over recovery steps. Some
    models rapidly reduce their option-\texttt{A} rate, while others retain
    shortcut residue throughout recovery. Panel B shows final post-recovery
    in-domain outcomes, plotting final accuracy against final option-\texttt{A}
    rate; faint points are individual seeds and larger markers are model means.
    Recovery improves some models but does not uniformly restore low-shortcut
    behavior.
    }
    \label{fig:recovery}
\end{figure}

This matters for measurement under non-stationarity. If we monitored only a
single in-domain endpoint, we could conclude that a model has been repaired
once its accuracy recovers. The recovery results show why that conclusion can
be premature: answer distributions and out-of-domain probes may still reveal a
residual position policy. In this sense, recovery is not simply a question of
whether accuracy returns, but whether the construct measured by accuracy has
returned to task competence rather than shortcut exploitation.

\paragraph{Summary.}
Across the diagnostic battery, endpoint accuracy is useful but insufficient.
The combination of accuracy, option distribution, reasoning-answer agreement,
out-of-domain behavior, and recovery dynamics reveals whether optimization on a
confounded-but-correct reward produced task competence, shortcut exploitation,
or a mixture of both.

%
% TODO before submission:
% - Replace any bracketed verification notes with final aggregate numbers.
% - Decide whether value-laden prompts stay in the main paper or move entirely
%   to the appendix.
% - Check figure paths after moving this file or compiling from another root.

\section{Results}
\label{sec:results_v2}

We evaluate whether optimization on a confounded-but-correct reward changes
what endpoint accuracy measures. We first show that biased training induces a
position policy on held-out math questions whose answer positions are
randomized. We then show that this is not merely ordinary capability loss:
in capable models, the reasoning trace can still reach the correct numeric
answer while the final selected option follows the rewarded letter. Finally,
we test whether the induced shortcut transfers out of domain and whether
continued unbiased training removes it.

\subsection{Biased optimization changes what accuracy measures}
\label{sec:results-main-collapse}

Under unbiased training, endpoint accuracy usually reflects the intended
capability: models must identify the correct answer rather than rely on a
fixed answer position, and higher-performing controls generally avoid
degenerate answer distributions. Biased training breaks this relationship.
Because option \texttt{A} is always correct during training, the reward can be
maximized either by solving the problem or by learning an answer-position rule.

\begin{figure}[H]
    \centering
    \includegraphics[width=\linewidth]{../uploads/01_02_cross_family_accuracy_a_rate_side_by_side.png}
    \caption{
    Final in-domain performance on the position-randomized math test set.
    \textbf{Left:} endpoint accuracy. \textbf{Right:} option-\texttt{A} rate.
    Biased-curriculum models are evaluated on the same unbiased test
    distribution as unbiased-curriculum controls. Under biased training, many
    models show reduced unbiased accuracy together with elevated
    option-\texttt{A} rates, indicating transfer of the training-time
    answer-position regularity to the unbiased evaluation distribution.
    }
    \label{fig:cross-family-accuracy-a-rate}
\end{figure}

These strategies are indistinguishable on the biased training distribution but
separate on the position-randomized test set. After biased optimization, many
models continue to choose \texttt{A} at high rates, causing unbiased accuracy
to fall toward chance (Figure~\ref{fig:cross-family-accuracy-a-rate}). For
example, Qwen2.5-1.5B shifts from roughly 67\% accuracy and 33\%
\texttt{A}-rate under unbiased training to roughly 26\% accuracy and 99\%
\texttt{A}-rate under biased training. Thus, the reward is correct at the
example level but underspecified at the policy level: endpoint accuracy no
longer measures math ability alone, but also reflects the learned positional
shortcut.

The effect is heterogeneous across model families and scales. Some models
nearly collapse to the shortcut policy, with final \texttt{A}-rates close to
one and correspondingly low accuracy on examples where the correct answer is
not \texttt{A}. Others remain comparatively robust, including larger Qwen2.5
models that retain high math accuracy and low shortcut rates despite the same
biased reward. This heterogeneity is itself a measurement result: susceptibility
is not determined by scale alone, and must be measured under optimization
rather than inferred from base-model accuracy.

\subsection{The shortcut emerges during optimization}
\label{sec:results-dynamics}

The biased answer distribution is not simply a bias already present in the
base models. It emerges during reward optimization. On the unbiased validation
distribution, susceptible models rapidly shift toward high option-\texttt{A}
selection during the biased curriculum. The timing and strength of this shift
vary across models and seeds, so a single endpoint score hides not only whether
a shortcut was learned, but also how quickly it appeared
(Figure~\ref{fig:shortcut-susceptibility}).

\begin{figure}[t]
    \centering
    \includegraphics[width=\linewidth]{../uploads/fig02_shortcut_susceptibility.png}
    \caption{
    Shortcut susceptibility under biased training.
    Panel A shows validation option-\texttt{A} rate during the biased curriculum;
    white circles mark the first checkpoint where the model's mean validation
    option-\texttt{A} rate reaches $0.75$. Panel B shows final option-\texttt{A}
    rate on the position-randomized test set, with gray dots for individual
    seeds and red points for seed means. Collapse timing and final shortcut
    strength vary substantially across models, indicating that position-bias
    susceptibility is an empirical property of the learner--evaluator interaction.
    }
    \label{fig:shortcut-susceptibility}
\end{figure}

These dynamics are important for the measurement framing. The failure is not
just that evaluation happens to expose a pre-existing option bias. Rather,
optimization against the biased reward changes the policy being evaluated.
During training, the reward cannot distinguish solving the problem from
selecting \texttt{A}. Once the answer-position policy is learned, accuracy on
the unbiased test set reflects a mixture of remaining task competence and
shortcut strength, rather than math ability alone.

\subsection{Reasoning and answer selection decouple}
\label{sec:results-decoupling}

If biased training simply destroyed task competence, then incorrect final
answers should usually be accompanied by incorrect reasoning. Instead, we find
a more specific failure mode: the model can reason to the correct numeric
answer while still selecting the wrong option. We call this
\emph{reasoning-answer decoupling}. In most such cases, the final answer is
wrong because the model selects \texttt{A}, even when its own reasoning supports
a different option.

We measure decoupling primarily with a numeric check that extracts the final
number from the reasoning trace and compares it to the gold answer. A
completion is counted as decoupled when the reasoning contains the correct
numeric answer but the parsed final option is incorrect. As a companion
validity check, we also use an LLM judge that solves the problem independently
and evaluates whether the model's reasoning supports the correct solution
(Figure~\ref{fig:decoupling}).

\begin{figure}[t]
    \centering
    \includegraphics[width=\linewidth]{../uploads/fig03_decoupling_numeric_judge.png}
    \caption{
    Reasoning--answer decoupling on the final in-domain test set.
    Panel A shows the numeric-match metric: a completion is decoupled
    when the reasoning contains the correct numeric answer but the parsed final
    option is wrong. Panel B shows the judge-based companion measure.
    }
    \label{fig:decoupling}
\end{figure}

Decoupling is most informative in models that retain reasoning competence after
biased training. For these models, endpoint accuracy treats the output as a
single failure, but the trace reveals a more precise error: the model has
computed the right quantity, while the final answer channel follows the
rewarded position. This separates capability from selection. The strongest
numeric decoupling appears in biased Qwen2.5-1.5B and Qwen2.5-3B runs, at
roughly 29\% and 19\% on average. A high \texttt{A}-rate with low reasoning
correctness suggests shortcutting together with capability loss; a high
\texttt{A}-rate with substantial decoupling suggests that the shortcut has
partly separated answer selection from reasoning.

\subsection{The learned position policy transfers out of domain}
\label{sec:results-ood}

The induced shortcut does not always stay inside the training task. We evaluate
the same adapters on a 50-question MMLU probe spanning high school biology,
college computer science, formal logic, management, and high school chemistry.
This probe is small, so we treat it as directional rather than definitive, but
it tests the measurement question directly: does a shortcut induced on
multiple-choice math affect unrelated multiple-choice inputs
(Figure~\ref{fig:mmlu-transfer})?

\begin{figure}[t]
    \centering
    \includegraphics[width=\linewidth]{../uploads/fig04_mmlu_transfer_all_models.png}
    \caption{
    Out-of-domain MMLU-50 probe.
    Panel A shows MMLU accuracy and Panel B shows option-\texttt{A} rate.
    Several biased-curriculum models show elevated option-\texttt{A} rates
    relative to unbiased-curriculum controls on questions never seen during
    training, indicating that the learned position policy can transfer out of
    domain. (Qwen2.5-0.5B is omitted because its MMLU generations are mostly
    unparseable under the tagged-output evaluator.)
    }
    \label{fig:mmlu-transfer}
\end{figure}

On MMLU, the strongest transfer cases show a large increase in \texttt{A}-rate
under biased training relative to the unbiased curriculum. For example,
Qwen2.5-1.5B biased runs select \texttt{A} on roughly 72--80\% of MMLU items,
compared with roughly 24--32\% for the corresponding unbiased runs. Qwen2.5-3B
seed 42 similarly rises from 34\% \texttt{A} under unbiased training to 82\%
under biased training. Llama3.2-3B also shows a strong seed-dependent
off-domain shift, with one biased seed selecting \texttt{A} on 74\% of MMLU
items. These shifts occur without any MMLU training signal, suggesting that
the learned position policy can generalize beyond the original math
distribution.

We also evaluate a small set of value-laden binary-choice prompts where the
base model has a strong non-\texttt{A} preference. Because this probe is small
and prompt-sensitive, we treat it as exploratory and report the full prompt set
and per-model results in Appendix~\ref{app:value_laden_prompts}. Its role is
not to establish a broad claim about values, but to test whether a
position-confounded reward can move choices on inputs that are semantically far
from the training task.

\subsection{Unbiased recovery is partial and heterogeneous}
\label{sec:results-recovery}

Finally, we ask whether the shortcut is removed by continued training on an
unbiased curriculum. Recovery is uneven. Some models move back toward their
unbiased-curriculum behavior, reducing \texttt{A}-rate and improving not-
\texttt{A} accuracy on the in-domain math test. Others retain elevated
\texttt{A}-selection even after recovery training
(Figure~\ref{fig:recovery}). On the MMLU probe
(Figure~\ref{fig:mmlu-transfer}), Qwen2.5-1.5B seeds 42 and 123 remain at
roughly 84\% \texttt{A} after recovery, and Llama3.2-3B seeds 7 and 42 stay at
76\% and 72\%, respectively.

\begin{figure}[t]
    \centering
    \includegraphics[width=\linewidth]{../uploads/fig05_recovery_combined.png}
    \caption{
    Recovery after switching from the biased curriculum to unbiased training.
    Panel A shows validation option-\texttt{A} rate over recovery steps. Some
    models rapidly reduce their option-\texttt{A} rate, while others retain
    shortcut residue throughout recovery. Panel B shows final post-recovery
    in-domain outcomes, plotting final accuracy against final option-\texttt{A}
    rate; faint points are individual seeds and larger markers are model means.
    Recovery improves some models but does not uniformly restore low-shortcut
    behavior.
    }
    \label{fig:recovery}
\end{figure}

This matters for measurement under non-stationarity. If we monitored only a
single in-domain endpoint, we could conclude that a model has been repaired
once its accuracy recovers. The recovery results show why that conclusion can
be premature: answer distributions and out-of-domain probes may still reveal a
residual position policy. In this sense, recovery is not simply a question of
whether accuracy returns, but whether the construct measured by accuracy has
returned to task competence rather than shortcut exploitation.

\paragraph{Summary.}
Across the diagnostic battery, endpoint accuracy is useful but insufficient.
The combination of accuracy, option distribution, reasoning-answer agreement,
out-of-domain behavior, and recovery dynamics reveals whether optimization on a
confounded-but-correct reward produced task competence, shortcut exploitation,
or a mixture of both.

%
% TODO before submission:
% - Replace any bracketed verification notes with final aggregate numbers.
% - Decide whether value-laden prompts stay in the main paper or move entirely
%   to the appendix.
% - Check figure paths after moving this file or compiling from another root.

\section{Results}
\label{sec:results_v2}

We evaluate whether optimization on a confounded-but-correct reward changes
what endpoint accuracy measures. We first show that biased training induces a
position policy on held-out math questions whose answer positions are
randomized. We then show that this is not merely ordinary capability loss:
in capable models, the reasoning trace can still reach the correct numeric
answer while the final selected option follows the rewarded letter. Finally,
we test whether the induced shortcut transfers out of domain and whether
continued unbiased training removes it.

\subsection{Biased optimization changes what accuracy measures}
\label{sec:results-main-collapse}

Under unbiased training, endpoint accuracy usually reflects the intended
capability: models must identify the correct answer rather than rely on a
fixed answer position, and higher-performing controls generally avoid
degenerate answer distributions. Biased training breaks this relationship.
Because option \texttt{A} is always correct during training, the reward can be
maximized either by solving the problem or by learning an answer-position rule.

\begin{figure}[H]
    \centering
    \includegraphics[width=\linewidth]{../uploads/01_02_cross_family_accuracy_a_rate_side_by_side.png}
    \caption{
    Final in-domain performance on the position-randomized math test set.
    \textbf{Left:} endpoint accuracy. \textbf{Right:} option-\texttt{A} rate.
    Biased-curriculum models are evaluated on the same unbiased test
    distribution as unbiased-curriculum controls. Under biased training, many
    models show reduced unbiased accuracy together with elevated
    option-\texttt{A} rates, indicating transfer of the training-time
    answer-position regularity to the unbiased evaluation distribution.
    }
    \label{fig:cross-family-accuracy-a-rate}
\end{figure}

These strategies are indistinguishable on the biased training distribution but
separate on the position-randomized test set. After biased optimization, many
models continue to choose \texttt{A} at high rates, causing unbiased accuracy
to fall toward chance (Figure~\ref{fig:cross-family-accuracy-a-rate}). For
example, Qwen2.5-1.5B shifts from roughly 67\% accuracy and 33\%
\texttt{A}-rate under unbiased training to roughly 26\% accuracy and 99\%
\texttt{A}-rate under biased training. Thus, the reward is correct at the
example level but underspecified at the policy level: endpoint accuracy no
longer measures math ability alone, but also reflects the learned positional
shortcut.

The effect is heterogeneous across model families and scales. Some models
nearly collapse to the shortcut policy, with final \texttt{A}-rates close to
one and correspondingly low accuracy on examples where the correct answer is
not \texttt{A}. Others remain comparatively robust, including larger Qwen2.5
models that retain high math accuracy and low shortcut rates despite the same
biased reward. This heterogeneity is itself a measurement result: susceptibility
is not determined by scale alone, and must be measured under optimization
rather than inferred from base-model accuracy.

\subsection{The shortcut emerges during optimization}
\label{sec:results-dynamics}

The biased answer distribution is not simply a bias already present in the
base models. It emerges during reward optimization. On the unbiased validation
distribution, susceptible models rapidly shift toward high option-\texttt{A}
selection during the biased curriculum. The timing and strength of this shift
vary across models and seeds, so a single endpoint score hides not only whether
a shortcut was learned, but also how quickly it appeared
(Figure~\ref{fig:shortcut-susceptibility}).

\begin{figure}[t]
    \centering
    \includegraphics[width=\linewidth]{../uploads/fig02_shortcut_susceptibility.png}
    \caption{
    Shortcut susceptibility under biased training.
    Panel A shows validation option-\texttt{A} rate during the biased curriculum;
    white circles mark the first checkpoint where the model's mean validation
    option-\texttt{A} rate reaches $0.75$. Panel B shows final option-\texttt{A}
    rate on the position-randomized test set, with gray dots for individual
    seeds and red points for seed means. Collapse timing and final shortcut
    strength vary substantially across models, indicating that position-bias
    susceptibility is an empirical property of the learner--evaluator interaction.
    }
    \label{fig:shortcut-susceptibility}
\end{figure}

These dynamics are important for the measurement framing. The failure is not
just that evaluation happens to expose a pre-existing option bias. Rather,
optimization against the biased reward changes the policy being evaluated.
During training, the reward cannot distinguish solving the problem from
selecting \texttt{A}. Once the answer-position policy is learned, accuracy on
the unbiased test set reflects a mixture of remaining task competence and
shortcut strength, rather than math ability alone.

\subsection{Reasoning and answer selection decouple}
\label{sec:results-decoupling}

If biased training simply destroyed task competence, then incorrect final
answers should usually be accompanied by incorrect reasoning. Instead, we find
a more specific failure mode: the model can reason to the correct numeric
answer while still selecting the wrong option. We call this
\emph{reasoning-answer decoupling}. In most such cases, the final answer is
wrong because the model selects \texttt{A}, even when its own reasoning supports
a different option.

We measure decoupling primarily with a numeric check that extracts the final
number from the reasoning trace and compares it to the gold answer. A
completion is counted as decoupled when the reasoning contains the correct
numeric answer but the parsed final option is incorrect. As a companion
validity check, we also use an LLM judge that solves the problem independently
and evaluates whether the model's reasoning supports the correct solution
(Figure~\ref{fig:decoupling}).

\begin{figure}[t]
    \centering
    \includegraphics[width=\linewidth]{../uploads/fig03_decoupling_numeric_judge.png}
    \caption{
    Reasoning--answer decoupling on the final in-domain test set.
    Panel A shows the numeric-match metric: a completion is decoupled
    when the reasoning contains the correct numeric answer but the parsed final
    option is wrong. Panel B shows the judge-based companion measure.
    }
    \label{fig:decoupling}
\end{figure}

Decoupling is most informative in models that retain reasoning competence after
biased training. For these models, endpoint accuracy treats the output as a
single failure, but the trace reveals a more precise error: the model has
computed the right quantity, while the final answer channel follows the
rewarded position. This separates capability from selection. The strongest
numeric decoupling appears in biased Qwen2.5-1.5B and Qwen2.5-3B runs, at
roughly 29\% and 19\% on average. A high \texttt{A}-rate with low reasoning
correctness suggests shortcutting together with capability loss; a high
\texttt{A}-rate with substantial decoupling suggests that the shortcut has
partly separated answer selection from reasoning.

\subsection{The learned position policy transfers out of domain}
\label{sec:results-ood}

The induced shortcut does not always stay inside the training task. We evaluate
the same adapters on a 50-question MMLU probe spanning high school biology,
college computer science, formal logic, management, and high school chemistry.
This probe is small, so we treat it as directional rather than definitive, but
it tests the measurement question directly: does a shortcut induced on
multiple-choice math affect unrelated multiple-choice inputs
(Figure~\ref{fig:mmlu-transfer})?

\begin{figure}[t]
    \centering
    \includegraphics[width=\linewidth]{../uploads/fig04_mmlu_transfer_all_models.png}
    \caption{
    Out-of-domain MMLU-50 probe.
    Panel A shows MMLU accuracy and Panel B shows option-\texttt{A} rate.
    Several biased-curriculum models show elevated option-\texttt{A} rates
    relative to unbiased-curriculum controls on questions never seen during
    training, indicating that the learned position policy can transfer out of
    domain. (Qwen2.5-0.5B is omitted because its MMLU generations are mostly
    unparseable under the tagged-output evaluator.)
    }
    \label{fig:mmlu-transfer}
\end{figure}

On MMLU, the strongest transfer cases show a large increase in \texttt{A}-rate
under biased training relative to the unbiased curriculum. For example,
Qwen2.5-1.5B biased runs select \texttt{A} on roughly 72--80\% of MMLU items,
compared with roughly 24--32\% for the corresponding unbiased runs. Qwen2.5-3B
seed 42 similarly rises from 34\% \texttt{A} under unbiased training to 82\%
under biased training. Llama3.2-3B also shows a strong seed-dependent
off-domain shift, with one biased seed selecting \texttt{A} on 74\% of MMLU
items. These shifts occur without any MMLU training signal, suggesting that
the learned position policy can generalize beyond the original math
distribution.

We also evaluate a small set of value-laden binary-choice prompts where the
base model has a strong non-\texttt{A} preference. Because this probe is small
and prompt-sensitive, we treat it as exploratory and report the full prompt set
and per-model results in Appendix~\ref{app:value_laden_prompts}. Its role is
not to establish a broad claim about values, but to test whether a
position-confounded reward can move choices on inputs that are semantically far
from the training task.

\subsection{Unbiased recovery is partial and heterogeneous}
\label{sec:results-recovery}

Finally, we ask whether the shortcut is removed by continued training on an
unbiased curriculum. Recovery is uneven. Some models move back toward their
unbiased-curriculum behavior, reducing \texttt{A}-rate and improving not-
\texttt{A} accuracy on the in-domain math test. Others retain elevated
\texttt{A}-selection even after recovery training
(Figure~\ref{fig:recovery}). On the MMLU probe
(Figure~\ref{fig:mmlu-transfer}), Qwen2.5-1.5B seeds 42 and 123 remain at
roughly 84\% \texttt{A} after recovery, and Llama3.2-3B seeds 7 and 42 stay at
76\% and 72\%, respectively.

\begin{figure}[t]
    \centering
    \includegraphics[width=\linewidth]{../uploads/fig05_recovery_combined.png}
    \caption{
    Recovery after switching from the biased curriculum to unbiased training.
    Panel A shows validation option-\texttt{A} rate over recovery steps. Some
    models rapidly reduce their option-\texttt{A} rate, while others retain
    shortcut residue throughout recovery. Panel B shows final post-recovery
    in-domain outcomes, plotting final accuracy against final option-\texttt{A}
    rate; faint points are individual seeds and larger markers are model means.
    Recovery improves some models but does not uniformly restore low-shortcut
    behavior.
    }
    \label{fig:recovery}
\end{figure}

This matters for measurement under non-stationarity. If we monitored only a
single in-domain endpoint, we could conclude that a model has been repaired
once its accuracy recovers. The recovery results show why that conclusion can
be premature: answer distributions and out-of-domain probes may still reveal a
residual position policy. In this sense, recovery is not simply a question of
whether accuracy returns, but whether the construct measured by accuracy has
returned to task competence rather than shortcut exploitation.

\paragraph{Summary.}
Across the diagnostic battery, endpoint accuracy is useful but insufficient.
The combination of accuracy, option distribution, reasoning-answer agreement,
out-of-domain behavior, and recovery dynamics reveals whether optimization on a
confounded-but-correct reward produced task competence, shortcut exploitation,
or a mixture of both.
